\MakeAbstract{
    八皇后问题规模不大，却很适合观察“建模方式、剪枝时机与实现细节”如何共同影响求解效率。基于这一点，本报告在同一套 Python 测试框架下实现并比较六类方法：命令式全排列后逻辑筛选、纯逻辑 \texttt{permuteq}、加入自定义不等关系后的纯逻辑求解、DFS、BFS 与位运算回溯。评估目标不只看是否得到 92 个解，也关注约束何时生效、每种方法的工程开销落在什么位置。

    实验过程中遇到的一个实际问题是：Kanren 0.2.3 缺少可直接调用的不等约束。为此，报告补写了 \texttt{neq\_rel} 与 \texttt{diagonal\_neq\_rel}，并将它们接入按列增量求解流程。这样做不改变问题语义，但能更早截断冲突分支，显著降低逻辑方法的无效搜索成本。

    在 Windows 11、Python 3.12.7、$8\times 8$ 棋盘、每种方法重复 3 次的设置下，六种实现均稳定返回 92 组解。平均耗时分别为：M1 1761.672 ms、M2 16507.190 ms、M3 2856.098 ms、DFS 6.003 ms、BFS 5.900 ms、位运算回溯 0.615 ms。结果表明：在逻辑编程路径中，约束前置是主要提速来源；在搜索规模接近时，位级状态表示能进一步压缩单节点处理时间。完整脚本与 JSON 结果已附后，便于复现实验并扩展到更一般的 $N$ 皇后规模。
}{八皇后问题，约束满足，逻辑编程，Kanren，图搜索，位运算优化}
